\subsection{General case}

In this section, G denotes a locally compact commutative group.

PROPOSITION 3. --- a) There exists an integer \(n \geqslant 0\) and a subgroup L such that G is the direct product of L and a subgroup isomorphic to \(\mathbf{R}^n\) , and moreover L admits a compact open subgroup K such that L/K is discrete;

\begin{enumerate}
\def\labelenumi{\alph{enumi})}
\setcounter{enumi}{1}
\tightlist
\item
  The group G is the union of an increasing filtered family of open subgroups, each being a projective limit of groups isomorphic to groups of the form \(R^{p} \times T^{q} \times Z^{r} \times A\) , where A is a finite group and p, q, r are positive integers.
\end{enumerate}

We prove \(b\) ). For every compact neighborhood V of \(e\) , let \(\mathrm{G_V}\) denote the subgroup of G generated by V. It is open, and by the corollary of prop. 2, the group \(\mathrm{G_V}\) is a projective limit of groups of the form \(\mathbf{R}^p\times \mathbf{T}^q\times \mathbf{Z}^r\times \mathrm{A}\) , where A is a compact group and \(p,q\) and \(r\) in \(\mathbf{N}\) . When V ranges over the compact neighborhoods of \(e\) , these subgroups \(\mathrm{G_V}\) form a filtered family (since \(G_{V}\) and \(G_{W}\) are contained in \(G_{V\cup W}\) for all compact neighborhoods V and W of e). Finally, the group G is the union of the subgroups \(G_{V}\) .

Let H be an open subgroup of G generated by a compact neighborhood of e. There exist a compact group K and positive integers p and q such that H is isomorphic to \(R^{p} \times Z^{q} \times K\) (prop. 1 of II, p.~246); identify H with this product. The canonical surjection of H onto the divisible group \(R^{p}\) extends to a morphism \(\pi\) of G onto \(R^{p}\) (lemma 3 of II, p.~245). It is a projector \(\pi\) of G onto \(R^{p}\), which is continuous since its restriction to the open subgroup H is continuous. Hence G is the direct product of \(R^{p}\) and of the kernel L of \(\pi\) (TG, III, p.~47, cor.). We have \(Z^{q} \times K = H \cap L\) , hence \(Z^{q} \times K\) is an open subgroup of L. Thus K is a compact open subgroup of L, and consequently L/K is discrete.

PROPOSITION 4. --- Let B \(_{G}\) be the set of elements of G that generate a relatively compact subgroup of G. Then B \(_{G}\) is a closed subgroup of G and B \(_{G}^{\perp}\) is the identity component of \(\widehat{G}\) .

The set \(B_{G}\) is a subgroup of G, since the product of two compact subsets of G is a compact subset of G.

Let H be an open subgroup of G generated by a compact neighborhood of e. There exist positive integers p and q and a compact group K such that the subgroup H is isomorphic to \(R^{p} \times Z^{q} \times K\) (prop. 1 of II, p.~246). If one identifies these groups, one sees that \(B_{G} \cap H = K\) is closed in H. Since the family of open subgroups H generated by compact neighborhoods is an open cover of G (for example, x belongs to the subgroup generated by \(U \cap \{x\}\) for every fixed compact neighborhood U of e), it follows that \(B_{G}\) is closed.

Now let us calculate \(B_{G}^{\perp}\). By prop. 3, a), there exists a positive integer \(n \geqslant 0\) and a group L admitting an open compact subgroup such that G can be identified with \(R^{n} \times L\) . Then \(B_{G}\) is identified with \(\{0\} \times B_{L}\) , and \(B_{G}^{\perp}\) with \(R^{n} \times B_{L}^{\perp}\). We are thus reduced to the case where G = L admits a compact open subgroup K.

We then have \(K \subset B_{G}\). If one identifies \(\widehat{G/K}\) with \(K^{\perp}\) (Theorem 4 of II, p.~226), the orthogonal \((\mathrm{B}_{\mathrm{G}}/\mathrm{K})^{\perp}\) of \(B_{G}/K\) in \(\widehat{G/K}\) is identified with \(B_{G}^{\perp}\). But on the other hand \(B_{G}/K = B_{G/K}\), and since \(K^{\perp}\) is an open subgroup of \(\widehat{G}\) (cor. 2 of II, p.~233), the identity component \(\widehat{G}_{0}\) of \(\widehat{G}\) is also the identity component of \(K^{\perp}\). Thus the assertion for the discrete group G/K is equivalent to that for G.

Finally, suppose that G is discrete. The group \(\widehat{G}\) is then compact (prop. 18 of II, p.~233). The neutral component \((\widehat{\mathrm{G}})_{0}\) is the intersection of the open subgroups of \(\widehat{G}\) (TG, III, p.~35, prop. 14) and a subgroup of \(\widehat{G}\) is open if and only if it is closed and of finite index, or equivalently if its orthogonal is finite (Corollary 2 of II, p.~233); Corollary 4 of II, p.~228 shows that \((\widehat{\mathrm{G}})_{0}^{\perp}\) is the union of the finite subgroups of G, which is none other than \(B_{G}\) since G is discrete. We conclude by duality.

COROLLARY 1. --- Suppose G is compact. Then the following conditions are equivalent:

\begin{enumerate}
\def\labelenumi{(\roman{enumi})}
\item
  The group G is connected;
\item
  The group \(\widehat{\mathbf{G}}\) is torsion-free;
\item
  The group G is divisible.
\end{enumerate}

This follows from prop. 4 and cor. 7 of II, p.~229, since \(B_{G} = G\) .

COROLLARY 2. --- Suppose G is compact. Then G is totally disconnected if and only if \(\widehat{\mathbf{G}}\) is a torsion group.

The group G is totally disconnected if and only if its identity component is reduced to \(\{e\}\) ; prop. 4 shows that this condition is equivalent to \(B_{\widehat{G}} = \widehat{G}\) . Since the group \(\widehat{G}\) is discrete (prop. 18 of II, p.~233), this amounts to saying that each element of \(\widehat{G}\) generates a finite group, hence that \(\widehat{G}\) is torsion.

COROLLARY 3. --- If G is connected, then G is divisible.

Indeed, there exist a compact connected group K and an integer \(n \geqslant 0\) such that G is isomorphic to \(R^{n} \times K\) (II, p.~248, prop. 3, where one must have L = K). Corollary 1 shows that K is divisible, and therefore G is divisible.

